% GENERATED FILE. Edit canonical lessons, then run npm run generate:books.
% canonical-source-hash: fnv1a64:8536c4fe336027ec
% canonical-lessons: MR-C513-karyasala, MR-C513-ades, MR-C513-vibhag, MR-C513-sadar-karne, MR-C513-sutrasancalan-karne

\chapter{Workshops and Orders}
\label{ch:mr-workshops-and-orders}
\hlchaptermodality{513}

\begin{chapteropening}
\textbf{By the end of this chapter:} \emph{I can talk about workshops, departments and presenting work.}

Complete the last lesson of chapter 513: I can talk about workshops, departments and presenting work.
\end{chapteropening}

\section[\texorpdfstring{\mr{कार्यशाळा} (kāryaśāḷā)}{कार्यशाळा (kāryaśāḷā)}]{\mr{कार्यशाळा} (kāryaśāḷā) — a workshop; a training session}

\label{lesson:MR-C513-karyasala}



\begin{quote}
Before the new one: say the Marathi for a supply, then the Marathi for to sell well.
\end{quote}

\subsection*{You'll want to know: \mr{कार्यशाळा}}
\textbf{\mr{कार्यशाळा}} — \emph{kāryaśāḷā} — \textquotedblleft{}a workshop; a training session\textquotedblright{}.

\begin{quote}
\mr{शिक्षकांसाठी} \mr{तीन} \mr{दिवसांची} \mr{कार्यशाळा} \mr{झाली}.
\end{quote}

\begin{quote}
\emph{śikṣakāṁsāṭhī tīn divsāñcī kāryaśāḷā jhālī.}
\end{quote}

\begin{quote}
There was a three-day workshop for teachers.
\end{quote}

Feminine.

\begin{cousinweb}
\textbf{\mr{कार्यशाळा}} (\emph{kāryaśāḷā}) is a tatsama compound of Sanskrit \emph{kārya}, \textquotedblleft{}work\textquotedblright{}, and \emph{śālā}, \textquotedblleft{}hall\textquotedblright{} (Monier-Williams).
\end{cousinweb}

\begin{grammarlens}[title={What you've built}]
One more word for describing what you lived and felt.
\end{grammarlens}

\subsection*{Your turn}
\begin{itemize}
\raggedright
  \item \emph{Say it:} \emph{kāryaśāḷā}
  \item \emph{Say it:} \emph{kāryaśāḷā}, once more
  \item \emph{Say it:} \emph{śikṣakāṁsāṭhī tīn divsāñcī kāryaśāḷā jhālī.} — There was a three-day workshop for teachers.
  \item \emph{Recall:} say the Marathi for a supply, then the Marathi for to sell well, then say \emph{kāryaśāḷā} again
\end{itemize}

\begin{tcolorbox}[breakable,colback=teal!4,colframe=teal!35!black,title={Before you move on}]
What does \mr{कार्यशाळा} mean? (\textquotedblleft{}A workshop; a training session\textquotedblright{}.) Say it once more.
\end{tcolorbox}
\section[\texorpdfstring{\mr{आदेश} (ādeś)}{आदेश (ādeś)}]{\mr{आदेश} (ādeś) — a command; an order}

\label{lesson:MR-C513-ades}



\begin{quote}
Before the new one: say the Marathi for to sell well, then the Marathi for a workshop.
\end{quote}

\subsection*{You'll want to know: \mr{आदेश}}
\textbf{\mr{आदेश}} — \emph{ādeś} — \textquotedblleft{}a command; an order\textquotedblright{}.

\begin{quote}
\mr{व्यवस्थापकाने} \mr{नवीन} \mr{आदेश} \mr{दिला}.
\end{quote}

\begin{quote}
\emph{vyavasthāpakāne navīn ādeś dilā.}
\end{quote}

\begin{quote}
The manager gave a new order.
\end{quote}

Masculine.

\mr{आदेश} \mr{देणे}, to order.

\begin{cousinweb}
\textbf{\mr{आदेश}} (\emph{ādeś}) is a tatsama, taken from Sanskrit \emph{ādeśa}, \textquotedblleft{}advice, instruction\textquotedblright{} (Monier-Williams).
\end{cousinweb}

\begin{grammarlens}[title={What you've built}]
One more word for describing what you lived and felt.
\end{grammarlens}

\subsection*{Your turn}
\begin{itemize}
\raggedright
  \item \emph{Say it:} \emph{ādeś}
  \item \emph{Say it:} \emph{ādeś}, once more
  \item \emph{Say it:} \emph{vyavasthāpakāne navīn ādeś dilā.} — The manager gave a new order.
  \item \emph{Recall:} say the Marathi for to sell well, then the Marathi for a workshop, then say \emph{ādeś} again
\end{itemize}

\begin{tcolorbox}[breakable,colback=teal!4,colframe=teal!35!black,title={Before you move on}]
What does \mr{आदेश} mean? (\textquotedblleft{}A command; an order\textquotedblright{}.) Say it once more.
\end{tcolorbox}
\section[\texorpdfstring{\mr{विभाग} (vibhāg)}{विभाग (vibhāg)}]{\mr{विभाग} (vibhāg) — a department; a division}

\label{lesson:MR-C513-vibhag}



\begin{quote}
Before the new one: say the Marathi for a workshop, then the Marathi for a command.
\end{quote}

\subsection*{You'll want to know: \mr{विभाग}}
\textbf{\mr{विभाग}} — \emph{vibhāg} — \textquotedblleft{}a department; a division\textquotedblright{}.

\begin{quote}
\mr{माझी} \mr{बहीण} \mr{हिशेब} \mr{विभागात} \mr{काम} \mr{करते}.
\end{quote}

\begin{quote}
\emph{mājhī bahīṇ hiśeb vibhāgāt kām karte.}
\end{quote}

\begin{quote}
My sister works in the accounts department.
\end{quote}

Masculine.

\begin{cousinweb}
\textbf{\mr{विभाग}} (\emph{vibhāg}) is a tatsama, taken from Sanskrit \emph{vibhāga}, \textquotedblleft{}distribution, a share\textquotedblright{} (Monier-Williams).
\end{cousinweb}

\begin{grammarlens}[title={What you've built}]
One more word for describing what you lived and felt.
\end{grammarlens}

\subsection*{Your turn}
\begin{itemize}
\raggedright
  \item \emph{Say it:} \emph{vibhāg}
  \item \emph{Say it:} \emph{vibhāg}, once more
  \item \emph{Say it:} \emph{mājhī bahīṇ hiśeb vibhāgāt kām karte.} — My sister works in the accounts department.
  \item \emph{Recall:} say the Marathi for a workshop, then the Marathi for a command, then say \emph{vibhāg} again
\end{itemize}

\begin{tcolorbox}[breakable,colback=teal!4,colframe=teal!35!black,title={Before you move on}]
What does \mr{विभाग} mean? (\textquotedblleft{}A department; a division\textquotedblright{}.) Say it once more.
\end{tcolorbox}
\section[\texorpdfstring{\mr{सादर} \mr{करणे} (sādar karṇe)}{सादर करणे (sādar karṇe)}]{\mr{सादर} \mr{करणे} (sādar karṇe) — to present; to submit}

\label{lesson:MR-C513-sadar-karne}



\begin{quote}
Before the new one: say the Marathi for a command, then the Marathi for a department.
\end{quote}

\subsection*{You'll want to know: \mr{सादर} \mr{करणे}}
\textbf{\mr{सादर} \mr{करणे}} — \emph{sādar karṇe} — \textquotedblleft{}to present; to submit\textquotedblright{}.

\begin{quote}
\mr{मीराने} \mr{बैठकीत} \mr{तिचा} \mr{अहवाल} \mr{सादर} \mr{केला}.
\end{quote}

\begin{quote}
\emph{mīrāne baiṭhkīt ticā ahvāl sādar kelā.}
\end{quote}

\begin{quote}
Meera presented her report at the meeting.
\end{quote}

Past, with -\mr{ने} (he …): \mr{त्याने} \mr{अहवाल} \mr{सादर} \mr{केला}, \emph{tyāne ahvāl sādar kelā}.

A report submitted or a song performed: \mr{गाणं} \mr{सादर} \mr{करणे}.

\begin{cousinweb}
\textbf{\mr{सादर} \mr{करणे}} (\emph{sādar karṇe}) is \mr{सादर}, a loan through Persian usually traced to Arabic \emph{ṣādir}, \textquotedblleft{}issued, going forth\textquotedblright{}, with \mr{करणे}.
\end{cousinweb}

\begin{grammarlens}[title={What you've built}]
One more word for describing what you lived and felt.
\end{grammarlens}

\subsection*{Your turn}
\begin{itemize}
\raggedright
  \item \emph{Say it:} \emph{sādar karṇe}
  \item \emph{Say it:} \emph{sādar karṇe}, once more
  \item \emph{Say it:} \emph{mīrāne baiṭhkīt ticā ahvāl sādar kelā.} — Meera presented her report at the meeting.
  \item \emph{Recall:} say the Marathi for a command, then the Marathi for a department, then say \emph{sādar karṇe} again
\end{itemize}

\begin{tcolorbox}[breakable,colback=teal!4,colframe=teal!35!black,title={Before you move on}]
What does \mr{सादर} \mr{करणे} mean? (\textquotedblleft{}To present; to submit\textquotedblright{}.) Say it once more.
\end{tcolorbox}
\section[\texorpdfstring{\mr{सूत्रसंचालन} \mr{करणे} (sūtrasañcālan karṇe)}{सूत्रसंचालन करणे (sūtrasañcālan karṇe)}]{\mr{सूत्रसंचालन} \mr{करणे} (sūtrasañcālan karṇe) — to host a programme}

\label{lesson:MR-C513-sutrasancalan-karne}



\begin{quote}
Before the new one: say the Marathi for a department, then the Marathi for to present.
\end{quote}

\subsection*{You'll want to know: \mr{सूत्रसंचालन} \mr{करणे}}
\textbf{\mr{सूत्रसंचालन} \mr{करणे}} — \emph{sūtrasañcālan karṇe} — \textquotedblleft{}to host a programme\textquotedblright{}.

\begin{quote}
\mr{शाळेच्या} \mr{कार्यक्रमाचं} \mr{सूत्रसंचालन} \mr{अरुणने} \mr{केलं}.
\end{quote}

\begin{quote}
\emph{śāḷecyā kāryakramācaṁ sūtrasañcālan aruṇne kelaṁ.}
\end{quote}

\begin{quote}
Arun hosted the school programme.
\end{quote}

Past, with -\mr{ने} (he …): \mr{त्याने} \mr{सूत्रसंचालन} \mr{केलं}, \emph{tyāne sūtrasañcālan kelaṁ}.

The event takes -\mr{चं}: \mr{कार्यक्रमाचं} \mr{सूत्रसंचालन} \mr{करणे}.

\begin{cousinweb}
\textbf{\mr{सूत्रसंचालन} \mr{करणे}} (\emph{sūtrasañcālan karṇe}) is a learned compound built from Sanskrit \emph{sūtra}, \textquotedblleft{}thread\textquotedblright{} (Monier-Williams), and \mr{संचालन}, \textquotedblleft{}running\textquotedblright{}: the host holds the thread of the programme.
\end{cousinweb}

\begin{grammarlens}[title={What you've built}]
One more word for describing what you lived and felt.
\end{grammarlens}

\subsection*{Your turn}
\begin{itemize}
\raggedright
  \item \emph{Say it:} \emph{sūtrasañcālan karṇe}
  \item \emph{Say it:} \emph{sūtrasañcālan karṇe}, once more
  \item \emph{Say it:} \emph{śāḷecyā kāryakramācaṁ sūtrasañcālan aruṇne kelaṁ.} — Arun hosted the school programme.
  \item \emph{Recall:} say the Marathi for a department, then the Marathi for to present, then say \emph{sūtrasañcālan karṇe} again
\end{itemize}

\begin{tcolorbox}[breakable,colback=teal!4,colframe=teal!35!black,title={Before you move on}]
What does \mr{सूत्रसंचालन} \mr{करणे} mean? (\textquotedblleft{}To host a programme\textquotedblright{}.) Say it once more.
\end{tcolorbox}
